\section{Incremental system evaluation}
\label{sec:incremental-evaluation}

The following evaluations extend the original architecture questions. They are
development measurements and local qualification, separate from the historical
cold experiment and from any future confirmatory population. Their profiles
and source versions are explicit; a partial combined controller is qualified
below, while the complete matched system comparison remains open.

\paragraph{Context substitution opportunity.}
We reconstructed all sixteen historical A/B request commitments across eight
unique families and checked the selected public source against its hashes.
All 254 selected windows were protected editable or test source; there were
zero eligible unchanged-dependency files and zero substituted windows.
The guarded substitution adapter therefore retained the raw source. Saved
native packets contained no explicit obligation CIDs and 14,471 unresolved
dependency-edge occurrences. Instruction text was retained, but this cohort
cannot establish explicit obligation preservation or dependency compression.

Using \texttt{tiktoken} 0.12.0 with \texttt{o200k\_base}, exact tokenization of
system/user message content totaled 91,783 for A, 107,870 for B, and 108,577 for
guarded substitution C: increases of 17.53\% and 18.30\% over A.
These counts exclude chat-template overhead and output, and are not Grok
provider usage. A second named encoding gives the same direction.
Eighteen constructed adapter controls separately check expansion and rejection;
they do not qualify a native producer. This result identifies a population
requirement: a compression experiment must include unchanged dependencies that
can actually be replaced under its stated policy.

\paragraph{Native dependency context and bounded projection.}
Three authored multi-module repositories supply six fresh native exact
dependency admissions. Twelve cold/incremental index and semantic-bundle root
comparisons agree; dependency edits invalidate prior capsules and unknown-call
mutations require raw fallback. Full transported native capsules total
8,399 serialized context-content
tokens versus 1,467 for raw source. A new restricted projection instead
transports source-bound integer equations. Six separate Z3 equivalence checks
and 252 concrete source executions check that projection; its 1,353
tokens are 7.77\% below raw. Plain AST source compaction totals
1,464 tokens. These are named-encoding offline counts, not provider savings.
The equation profile excludes floats, custom objects, side effects and unsupported
syntax. The compact projection is an evaluated addition; the full native capsule
payload itself is larger than raw source in these examples.

\paragraph{Repository context and preparation.}
Native repository custody, current-source expansion and invalidation checks
pass on bounded XML/pyparsing scopes and on both full declared Requests
inventories (97 files each); these profiles permit zero raw-source omissions.
Five full pyparsing producer processes reach their 600-second deadline. A
subsequent immutable-state identity profile completes full preparation in
566.254 seconds, followed by a passing independent native consumer; all
302 XML/pyparsing regular files are retained and 22 negative controls pass.
This establishes completion under the bound in one shared-host run. Preparation
is still substantial, and these profiles establish no compression or replicated
speedup. A subsequent full comment-change comparison exposes 1,568 differing
graph records. A rescanning guard restores byte equality but fails a later
import-inventory case. Preserving original frontend observations then gives
exact full equality while reusing 153 of 154 Python files. Incremental takes
607.971 seconds versus 440.276 cold, adds a 53.75 MB sidecar and incurs separately
metered bootstrap costs. Useful incremental economics remain open.
Appendix~\ref{sec:incremental-context-preparation} retains source scopes,
identity checks, earlier failures and actual costs.

\paragraph{Live context substitution.}
The three dependency-bearing development repositories also support a matched
raw-versus-equation repair pilot. Both arms pass all three separate cold
behavioral checks; all six requests return usage. Provider input totals are
4,014 raw versus 3,817 compact (4.91\% fewer),
and total usage is 4,866 versus 4,637 (4.71\%
fewer). Both arms use three requests with the same Grok 4.6 low-effort profile
and resource ceilings as the routing pilot. Cached-input totals differ, so
charge differences cannot be attributed solely to compaction. This tests a
restricted, verified representation on trivial constructed repairs; dynamic
source expansion and broader repository quality remain unmeasured. The result
is separate from, and must not be added to, the routing reduction.

\paragraph{Checked symbolic closure and live residuals.}
A new development caller invokes the native Python analyzer, bounded CEGIS
repair, visible behavioral checks and native routing. Four constructed tasks
include false success without observation, swallowed backend failure, interval
clamping and stable deduplication. Both arms receive the same source, objective
and public checks; separate Docker execution scores broader behavioral checks.
The first two defects fit the declared deterministic grammar. The remaining two
require the same model residual in both arms. Internal CEGIS marker gates are
not counted as theorem proofs; behavioral checks justify only these bounded
obligations.

\begin{table}[htbp]
\centering
\caption{Constructed four-task routing development comparison. One observation
per task and arm; no population inference.}
\label{tab:routing-development}
\begin{tabular}{lrrr}
\toprule
Arm & Independently passed & Model calls & Total provider tokens \\
\midrule
Model first & 4/4 & 4 & 4,865 \\
Symbolic first plus residual & 4/4 & 2 & 2,430 \\
\bottomrule
\end{tabular}
\end{table}

All six Grok 4.6 requests returned usage, under the same low-effort profile,
\texttt{max\_tokens=2048} request parameter and 240-second deadline, without retries.
The observed reduction is 50.05\% on this suite. Eliminating two requests
accounts for 2,282 baseline tokens; response variation on the residual tasks
accounts for the remaining 153-token net difference. Provider cache state was
not controlled. There were no auxiliary workflow model calls or model-assisted
setup in this caller; experiment authoring and implementation effort are outside
that runtime denominator. The initial expired-credential attempt dispatched
zero requests and remains a separate development record. General repository
performance and full-system setup amortization remain unmeasured.

\paragraph{Supplementary scorer validation.}
A post-development audit exposes five false admissions among six crafted
incorrect candidates in the original pilot scorers. Candidate interference
with oracle helpers or premature process termination could masquerade as
passing execution. A new isolated scorer separates the oracle from the candidate
worker, restricts public dependency imports and requires observed calls and
positive completion. It rejects all six crafted candidates and all seven
unrepaired controls. All ten unique saved candidate bodies still pass, preserving
the eight routing and six context record outcomes above. This supplementary
rescore supports these saved outcomes; the small control set does not establish
general scorer robustness. Original scorers and observations remain retained.

\paragraph{Lifecycle-preserving test reuse.}
A proposed local evaluation adapter uses native cache admission and Ed25519
runner attestations issued only after actual passing setup/call/teardown.
Its explicitly enrolled pure workload and immutable fixture profile support
valid call-phase reuse while setup and teardown still execute. Fourteen
prespecified cases cover unchanged work, source/test/fixture/configuration/
environment mutations, missing or tampered evidence, revocation, unsupported
fixture values and lifecycle failure. All fourteen match separate cold-process
outcomes. Reuse followed by failing teardown correctly fails the whole item;
the result is not zero-error evidence for arbitrary pytest programs.

Five alternating warm/cold timing pairs omit all five warm test-body executions
while retaining five setups and teardowns in each arm. Total host wall time was
1.281 seconds cold and 1.240 seconds warm, excluding the separately recorded
prime and issuance costs. This small descriptive difference does not establish
a reliable net speedup or token savings. The default production plugin probe
executed the body in both write and read/write modes and produced no positive
reuse. Thus the result qualifies the proposed bounded adapter, not deployment
of the native default reuse path. Earlier failed development invocations and
the historical false-reuse observation remain retained.
The measured 0.256-second cold prime and 0.006-second issuance exceed
the observed wall-time saving over these five repetitions; a net benefit at
that horizon is not demonstrated.

\paragraph{Persistent work across process termination.}
Three actual processes exercise the native session journal/checkpoint with
the same checked deterministic caller. Duplicate pending notifications coalesce.
The first process exits abruptly after one completion; the next resumes exactly
one unfinished scan, and a third restart finds no pending work. Each actual
execution passes the independent behavioral oracle, and no original attempt
executes twice. However, a fresh notification for a completed snapshot schedules
and executes analysis again. Journal continuity is supported by this case;
completed-result memoization, accepted-root publication and cross-task token
savings are not established.

\paragraph{Publication and recovery qualification.}
Forty-seven existing source-shipped tests pass in a fresh compatible native
profile. They cover immutable-root publication, stale writers, competing CAS
updates, seal storage and restart recovery at injected persistence boundaries.
These are constructed software qualifications, including placeholder proof
records and exception-based interruptions. They do not establish a newly
integrated evidence-producing workflow, actual process-kill tolerance at each
boundary, or measured parallel-sealing savings.

\paragraph{Evidence-backed sealing and completed recovery.}
A new caller independently replays finite computation receipts, checks exact
source/input/policy bindings and complete coverage, then invokes native
integrity-only seal verification and publication. A full checkpoint and two
delta seals become successive current roots; a competing delta loses stale
parent CAS. Four prepared states cover 700 independently checked state/input
observations, using five computation subprocesses and three unchanged-receipt
reuses. Five constructed mutations reject. Native digest verification alone
accepts three internally consistent but semantically wrong or stale receipts;
the caller's independent replay and bindings are therefore substantive gates.

Two real process terminations preserve the expected accepted roots. Native
recovery recognizes a valid pre-CAS seal but leaves its transition pending;
retrying the same publication is refused. A separately evaluated recovery
caller rechecks the persisted evidence and completes native CAS/WAL phases
with zero new computation calls. A subsequent fresh process consumes that
accepted root and publishes the next changed state, reaching a four-seal chain.
This establishes bounded continuation under an explicit added caller. It does
not establish native automatic recovery, theorem proofs, power-loss tolerance
or a compute saving: the independent receipt checker still runs on reuse.
Four paired trials separately compare actual preparation of two cold operations:
all eight preparations pass with identical receipt and seal content. Median
preparation time is 63.21 ms sequential versus 32.10 ms parallel;
maximum sampled aggregate process RSS rises from 44.66 to 55.54 MiB.
The host is not exclusive, sampled RSS counts shared pages per process, and
these tiny preparation trials measure neither publication latency nor a
whole-system speedup.

\paragraph{Procedure synthesis and stored candidates.}
A bounded native synthesis/body-store comparison completes seven of eight units
in both cold and stored arms: 3,883 correct fresh-query decisions, zero false
admissions and 554 abstentions on unsupported parity. Stored candidates require
six synthesis calls and 23 tried candidates versus eight and 29 cold, under exact
source-catalog and policy bindings. A separately frozen fixed-template control
matches quality with no synthesis, so retrieval has no demonstrated advantage
over that simpler control. Fresh-process replay, binding mutations and scaffold
sensitivity are checked; this initial executor/store profile has no model calls.
A later signed native TEST integration executes the exact saved candidate on
136 new bounded inputs through all six native handlers. All decisions match
the oracle; 30 evidence negatives and 13 runtime controls pass. Its 17.243-second
worker cost plus preparation does not establish savings. Appendix~\ref{sec:signed-native-procedure}
separates this tested execution from repository applicability and promotion.

\paragraph{World-state memory and actual subsequent actions.}
A native typed transition producer, persistent block/SQLite store and admission
path are connected by a new evaluation dispatcher. Eleven tasks per arm run in
fresh interpreters against a closed three-action integer service; selected
actions always execute and their current outputs are independently checked.
All 44 task/arm outcomes pass. Native memory uses 16 action executions versus
27 without memory, but exact JSON retrieval and complete-condition templates
also use 16. Charging each arm the common three-action seed gives 19 versus
30. Changed state, source and declared environment invalidate prior observations;
fallback creates fresh evidence consumed by the next process.
Native total process time is 8.947 seconds versus 8.739 without
memory, including common import/storage work and excluding the separately
recorded shared seed. Thus fewer search attempts do not establish a wall-time
advantage. The frozen dispatcher uses retrieved typed records directly: an
audit finds that the native compiler emits incorrect record identifiers in
all eleven compiled proposals, so compiler-mediated dispatch is not qualified
by this run. These results establish bounded storage/lookup consumption, not
learned world-model benefit, transfer, production integration or token savings.
A separately frozen compiler correction selects each record's own identity.
The ensuing consumer resolves compiler-selected records from the native store
before actual execution: the seed and eleven later cases pass, using 19 actions
and 25 fresh service processes. Changed source, state and environment still
invalidate prior observations. Independent checks preserve all original
comparison artifacts. This qualifies the corrected compiler-mediated path;
it does not change the original simple-retrieval/template parity or supply
learned transfer and token-savings evidence.

\paragraph{Native compositional reasoning.}
An existing original-source vertical slice was run with its compatible nested
datasets tree. A baseline component contract discharged, an injected fault was
refuted, and repair restored discharge. Localization produced an UNSAT core and
a cvc5 interpolant separately checked by Z3; the final artifact was verified by
semantic rescanning, composition redischarge and test replay. Selected and full
tests passed. Host elapsed time was 19.309 seconds. This is one existing
constructed fixture with a prescribed fault and expected repair, not a held-out
synthesis experiment. Internal benchmark estimates are not counted as measured
token or compute savings. An earlier source profile failed to import the
required composition module and is preserved separately. Solver validation
does not imply kernel reconstruction or all-language correctness.

\paragraph{Native proof production and statement binding.}
The source-built native backend produces three v3 Merkle/count proofs and three
honest v2 implication-chain proofs. For v3, primary setup takes 7.810
seconds, proving 12.048--12.474 seconds and verification
3.053--5.147 milliseconds. Tested malformed witnesses, proofs and
keys reject. Both profiles require exact outer-statement checks at the caller:
the raw verifier can accept changed displayed metadata when its scalar public
inputs remain unchanged. These are measured bounded circuit relations, not
proofs that arbitrary tests executed or production trust admission. Stronger
v2 textual-binding assurance remains unconfirmed, and the v5 pytest-attestation
profile is unsupported. Appendix~\ref{sec:incremental-native-proof} retains the
separate profiles, mutation observations and v2 costs. A subsequent native
archive producer/consumer binds four actual accepted controller events to one
honest v3 proof, with seven mutation rejections. A second integration consumes
the actual accepted Requests repair, restores all 97 source files and rejects
eight mutations. Its 11.979-second proof and 5.256-millisecond verification
are post-cohort overhead; original receipts still supply execution evidence.

\paragraph{Combined controller and native publication.}
Constructed thirteen-event reasoning and eight-event publication profiles join
the actual native services and exercise real interruptions, native full/delta
publication and fresh accepted-source consumption. A separate Requests curator
fixture measures 377.226 seconds of phases including full context
preparation. These zero-provider controls establish bounded continuation,
not token savings. Appendix~\ref{sec:incremental-controller-publication} retains
their complete costs, failures and unsigned integrity-only trust boundary.

\paragraph{Natural two-issue comparisons.}
Two frozen Requests 5414-to-6028 development cohorts compare four arms with
equal public access and four-call trajectory ceilings. Every arm in each
cohort accepts the first issue and fails the second. First requests within
each cohort are byte-identical; B permits no source omission, while C/D
establish no successful symbolic closure. Partial D publishes and retrieves
accepted source. The first cohort's A/C/D continuations exceed their source
bound; B spends its remaining calls on failed candidates. A new conflict-first
retriever keeps required source within the same bounds in the second cohort.
All sixteen calls then execute, including expansion requests and three outputs
rejected by a file-mode transport defect. Table~\ref{tab:requests-cold} retains
both cohorts separately. In the second, B/C/D use 10.43\%/12.01\%/7.77\% more
tokens than A at the same observed acceptance count. Uncontrolled server caching,
one exposed family and adaptive follow-up preclude causal component or general
efficiency claims. Complete-trajectory cost ratios remain undefined.

\begin{table}[htbp]
\centering
\caption{Natural Requests cold development cohorts, retained separately.
Every arm accepts 1/2 nodes and completes 0/1 trajectories in each cohort.
Second-cohort wall includes nested services; first-cohort timing limits remain
in the appendix.}
\label{tab:requests-cold}
\begin{tabular}{lrrrrr}
\toprule
Arm & First calls & First tokens & Second calls & Second tokens & Second seconds \\
\midrule
A & 1 & 14,428 & 4 & 56,543 & 227.882 \\
B & 4 & 65,695 & 4 & 62,439 & 249.385 \\
C & 1 & 14,589 & 4 & 63,336 & 415.743 \\
D (partial) & 1 & 14,269 & 4 & 60,937 & 489.516 \\
\bottomrule
\end{tabular}
\end{table}

\paragraph{Remaining system questions.}
These increments establish local mechanisms and expose limitations. They leave
open matched context-substitution efficacy, persistent multi-task quality/cost,
procedure and world-memory utility on real tasks, production reuse, governed
publication, integrated native proof admission, and a held-out comparison with complete
startup, fallback, failed-attempt and recovery accounting.
Additional benchmark recruitment and environment qualification are reported in
Appendix~\ref{sec:incremental-recruitment}; they add no model-quality outcomes.
